\section{目标与源项定义}

本章计算目的：明确本报告的使用边界。本报告给出典型水下目标电场源项的量级预测、模型可比性和参数包络，不根据强度或模型可信度替用户给出实验取舍。

\subsection{报告目标}

本报告选取典型潜艇/水下目标场景，对 100、300、500、1000 m 距离尺度下的目标相关电场量级进行统一整理。输出结果采用 min/reference/max 三层格式，其中 min 和 max 来自参数包络，reference 来自中等参考场景。所有结果只表示模型估算量级，不代表任何真实平台。变频器--水泵系统高频电磁背景单独作为当前实验背景场处理，不进入目标源项 100--1000 m 总表。

\subsection{源项定义}

本报告计算和标注以下四类目标相关源项：

\begin{enumerate}
  \item UEP/ICCP 静态或慢变电场：由腐蚀、防腐和外加电流系统共同形成的等效静态或慢变电偶极签名。
  \item 轴频基频：等效偶极矩随轴频产生的基频调制分量。
  \item 轴频二/三阶谐波：轴频调制的二阶和三阶线谱分量。
  \item 尾流/流动诱导电场：尾流中心线局部流速与地磁场耦合产生的感生电场估算。
\end{enumerate}

另设一类实验背景源：变频器--水泵系统高频电磁背景。该源表示小型变频器 VFD 接入电源并驱动 2.2 kW 三相水泵时，由 VFD PWM 开关、电机电缆、泵体、水体、接地路径、DAQ 和传感器链路耦合形成的实验背景场，不作为真实水下目标源项远场预测。

\subsection{中立性约束}

最终表和结果读取说明保持客观：可以比较幅值范围、频率特征和模型可信度，但不根据这些比较给出实验取舍。尾流模型可信度可以标注为低于 UEP/轴频远场模型；轴频基频可以标注为幅值低于静态场但具有稳定频率线谱；变频器--水泵系统高频电磁背景只按频率相关共模/漏电流等效偶极模型给出，用于当前实验背景量级对照。

本章对客观参考表的作用：限定全文只输出量级、频率特征、距离含义和模型边界，不输出实验排序或取舍结论。
